\section{Cost Allocation Tags}

\begin{itemize}
    \item With Tags we can track resources that relate to each other
    \item With Cost Allocation Tags we can enable detailed costing reports
    \item Just like Tags, but they show up as columns in Reports
    \item AWS Generated Cost Allocation Tags
    \begin{itemize}
        \item Automatically applied to the resource you create
        \item Start with Prefix aws: (e.g. aws: createdBy)
        \item They're not applied to resources created before the activation
    \end{itemize}
    \item User tags
    \begin{itemize}
        \item Defined by the user
        \item Start with Prefix user:
    \end{itemize}
    \item Cost Allocation Tags just appear in the Billing Console
    \item Takes up to 24 hours for the tags to show up in the report
\end{itemize}


\section{AWS Tag Editor}
\begin{itemize}
    \item Allows you to managed tags of multiple resources at once
    \item You can add/update/delete tags
    \item Search tagged/untagged resources in all AWS Regions
\end{itemize}


\section{Trusted Advisor}
\begin{itemize}
    \item No need to install anything - high level AWS account assessment
    \item Analyse your AWS accounts and provides recommendation
    \item Cost Optimisation
    \item Performance
    \item Security
    \item Fault Tolerance
    \item Service Limits
    \item Operational Excellence
    - Core Checks and recommendations - all customers
    \item Can enable weekly email notification from the console
    \item Full Trusted Advisor - Available for Business and Enterprise support plans
    \item Ability to set CloudWatch alarms when reaching limits
    \item Programmatic Access using AWS support API
\end{itemize}

\begin{table}[h!]
    \centering
    \begin{tabular}{|c|c|c|c|c|}
        \hline
        Column1                                  & Basic Support                                                        & Column3 & Column4 & Column5 \\ \hline
        AWS Trusted Advisor Best Practice Checks & 7 Core Checks                                                        & Data3   & Data4   & Data5   \\ \hline
        Enhanced Technical Support               & 24x7 customer service, documentation, whitepapers and support forums  & Data8   & Data9   & Data10  \\ \hline
        Case Severity / Response Times           & Data12                                                               & Data13  & Data14  & Data15  \\ \hline
        Data16                                   & Data17                                                               & Data18  & Data19  & Data20  \\ \hline
    \end{tabular}
    \caption{Example Table with 4 Rows and 5 Columns}
    \label{tab:example_table}
\end{table}

\subsubsection{Trusted Advisor - Good to Know}
\begin{itemize}
    \item Can check if an S3 bucket is made public
    \item But cannot check for S3 objects that are public inside of your bucket
    \item Use Amazon EventBridge / S3 Events instead / AWS Config Rules

    \item Service Limits
    \begin{itemize}
        \item Limits can only be monitored in Trusted Advisor (cannot be changed)
        \item Cases must be created manually in AWS Support Centre to increase limits
        \item OR use the AWS Service quotas service
    \end{itemize}
\end{itemize}


\section{AWS Service Quotas}
\begin{itemize}
    \item Notify you when you're close to a service quota value threshold
    \item Create CloudWatch Alarms on the Service Quotas console
    \item Example: Lambda concurrent executions
    \item Helps you know if you need to request a quota increase or shutdown resources before limit is reached
\end{itemize}
%Diagram


\section{EC2 Launch Types and Savings Plans}
\begin{itemize}
    \item On Demand Instances - short workload, predictable pricing, reliable.
    \item Spot Instances - short workloads for check, can lose instances (not reliable)
    \item Reserved: (Minimum 1 year)
    \item Reserved Instances - long workloads
    \item Convertible Reserved Instances - long workloads with flexible instances
    \item Dedicated Instances: no other customers will share you hardware
    \item Dedicated Hosts: book an entire physical server, control instance placement
    \item Great for software licenses that operate at the core, or socket level
    \item Can define host affinity so that instance reboots are kept on the same host
\end{itemize}

\subsection{AWS Savings Plan}
\begin{itemize}
    \item New pricing model to get a discount based on long-term usage
    \item Commit to a certain type of usage: ex \$10 per hour for 1 to 3 years
    \item Any usage beyond the savings plan is billed at the on-demand price
\end{itemize}

\begin{itemize}
    \item EC2 Instance Savings plan (72\% - same discount as Standard RIs)
    \item Select instance family and locked to a specific region
    \item Flexible across size, OS (Windows to Linux) tenancy. (dedicated or default)
    \item Compute Savings Plan
    \item Ability to move between instance family, region, compute type and OS and tenancy
    \item SageMaker Savings plan (up to 64\% off)
\end{itemize}


\section{S3 Cost Savings}

\subsection{S3 Storage Classes}
\begin{itemize}
    \item Amazon S3 Standard - General Purpose
    \item Amazon S3 Standard-Infrequent Access (IA)
    \item Amazon S3 One Zone-Infrequent Access
    \item Amazon S3 Glacier Instant Retrieval
    \item Amazon S3 Glacier Flexible Retrieval
    \item Amazon S3 Glacier Deep Archive
    \item Amazon S3 Intelligent Tiering
\end{itemize}

Can move between classes manually or using S3 lifecycle configurations

\subsection{S3 - Other Cost Savings}
\begin{itemize}
    \item S3 Lifecycle Rules: transition objects between tiers
    \item Compress Objects - to save space
    \item S3 Requester Pays:
    \item In general, bucket owners pay for all Amazon S3 storage and data transfer costs associated with their bucket
    \item With Requester Pays buckets, the requester instead of the bucket owner pays the cost of the request and the data downloaded from the bucket
    \item The bucket owner always pays the cost of storing data
    \item Helpful when you want to share large datasets with other accounts
    \item If an IAM role is assumed the owner account of that role pays for the request
\end{itemize}


\section{S3 Storage Classes - Reminder}

\subsection{S3 Storage Classes}
\begin{itemize}
    \begin{itemize}
        \item Amazon S3 Standard - General Purpose
        \item Amazon S3 Standard-Infrequent Access (IA)
        \item Amazon S3 One Zone-Infrequent Access
        \item Amazon S3 Glacier Instant Retrieval
        \item Amazon S3 Glacier Flexible Retrieval
        \item Amazon S3 Glacier Deep Archive
        \item Amazon S3 Intelligent Tiering
    \end{itemize}

    Can move between classes manually or using S3 lifecycle configurations

    \subsection{S3 Durability and Availalbility}
    - Durability:
    \begin{itemize}
        \item High Durability (99.9999999999, 11 nines) of objects across multiple AZ
        \item If you store 10,000,000 object with Amazon S3, you can on average expect to incur a loss of a single object once every 10,000 years (nice....but basic math)
        \item Same for all storage classes
    \end{itemize}

    - Availability
    \begin{itemize}
        \item Measures how readily available a service is
        \item Varies depending on storage class
        \item Example: S3 standard has 99.99\% availability = not available 53 minutes a year
    \end{itemize}

    S3 Standard - General Purpose
    \begin{itemize}
        \item 99.99\% Availability
        \item Used for frequently accessed data
        \item Low latency and high throughput
        \item Sustain 2 concurrent facility faiures
        \item Use Cases: Big Data analytics, mobile and gaming applications, content and distribution
    \end{itemize}

    S3 Storage Classes - Infrequent Access
    \begin{itemize}
        \item For data that is less frequently accessed, but requires rapid access when needed
        \item Lower cost than S3 standard
    \end{itemize}


    \begin{itemize}
        Amazon S3 Standard-Infrequent Access (S3 Standard-IA)
        \item 99.9\% Availability
        \item Use cases: Disaster Recovery, backups
    \end{itemize}

    Amazon S3 One Zone-Infrequent Access (S3 One Zone-IA)
    \begin{itemize}
        \item High durability (99.999999999) in a single AZ; data lost when AZ is destroyed
        \item 99.5\% Availability
        \item Use Cases: Storing secondary backup copies of on-premise data or data you can recreate
    \end{itemize}

    Amazon S3 Glacier Storage Classes
    \begin{itemize}
        \item Low-cost object storage meant for archiving / backup
        \item Pricing: price for storage + object retrieval cost
    \end{itemize}

    - Amazon S3 Glacier Instant Retrieval
    \begin{itemize}
        \item Millisecond retrieval, great for data accessed once a quarter
        \item Minimum storage duration of 90 days
        \item Amazon S3 Glacier Flexible Retrieval (Formerly Amazon S3 Glacier)
        \item Expedited (1 to 5 minutes), Standard (3 to 5 hours), Bulk (5 to 12 hours) - free
        \item Minimum storage duration of 90 days
        \item Amazon S3 Glacier Deep Archive - for long term storage
        \item Standard (12 hours), Bulk (48 Hours)
        \item Minimum Storage duration of 180 days
    \end{itemize}
\end{itemize}

\subsection{S3 Intelligent Tiering}
\begin{itemize}
    \item Small monthly monitoring and auto-tiering fee
    \item Moves objects automatically between Access Tiers based on usage
    - There are no retrieval charges in S3 Intelligent - Tiering
\end{itemize}

- Frequent Access tier (automatic): default tier
\begin{itemize}
    \item Infrequent Access Tier (automatic): objects not accessed for 30 days
    \item Archive Instant Access tier (automatic): objects not accessed for 90 days
    \item Archive Access tier (optional): configurable from 90 days to 700+ days
    \item Deep Archive Access tier (optional): config from 180 days to 700+ days
\end{itemize}

\subsubsection{S3 Storage Classes Comparision}
%classic table here

\subsubsection{S3 Storage Classes - Price Comparison}
%Examples


\section{AWS Budgets and Cost Explorer}

\subsection{AWS Budgets}
\begin{itemize}
    \item Create budget and send alarms when costs exceeds the budget
    \item 4 Types of budgets: Usage, Cost, Reservation, Savings Plans
    \item For Reserved Instances (RI)
    \item Track utilisation
    \item Supports EC2, ElastiCache, RDS, Redshift
    \item Up to 5 SNS notifications per budget
    \item Can filter by: Service, Linkedin Account, Tag, Purchase Option, Instance Type, Region, Availability Zone, API Operations, etc......
    \item Same options as AWS Cost Explorer
    \item 2 budgets are free than \$0.002 / day / budget
\end{itemize}

\subsection{Budget Actions}
\begin{itemize}
    \item Run actions on your behalf when a budget exceeds a certain cost or usage threshold
    \item Supports 3 actions types
    \item Applying an IAM Policy to a user, group or IAM role
    \item Applying Service Control Policy (SCP) to an OU
    \item Stop EC2 or RDS Instances
    \item Actions can be executed automatically or require a workflow approval process
    \item Reduced unintentional overspending in your account.
\end{itemize}
%Add diagram

\subsection{Centralised Budget Management}
%Diagram

\subsection{DeCentralised Budget Management}
%Diagram

\subsection{Cost Explorer}
\begin{itemize}
    \item Visualise, understand and manage your AWS costs and usage over time
    \item Create custom reports that analyse cost and usage data
    \item Analyse you data at a high level: total costs and usage across all accounts
    \item Or Monthly, hourly, resource level granularity
    \item Choose an optimal Savings PLan (to lower prices on your bill)
    \item Forecast usage up to 12 months based on previous usage
\end{itemize}

%Diagram Cost- Explorer example


\section{AWS Compute Optimiser}
\begin{itemize}
    \item Reduce costs and improve performance by recommending optimal AWS resources for you workloads
    \item Helps you choose optimal configurations and right-size your workloads (over/under provisioned)
    \item Uses Machine Learning to analyse your resources configurations and their utilisation CloudWatch metrics
    \item Supported resources
    \item EC2 Instances
    \item EC2 Auto Scaling Groups
    \item EBS volumes
    \item Lambda functions
    \item Lower your costs by up to 25\%
    \item Recommendations can be exported to S3
\end{itemize}

\subsection{Computer Optimiser - CloudWatch Agent}
\begin{itemize}
    \item Needed to analyse Memory Utilisation
    \item Not needed for CPU, NetworkIn/Out, DiskReadOps, DiskWriteOps
\end{itemize}
%Diagram


\section{EC2 Reserved Instance}
\begin{itemize}
    \item Reserved Instances in an AWS Organisation
    \item All accounts share the Reserved Instances and Savings Plan
    \item The payer account (Management account) of an organisation can turn off Reserved Instance (RI) discount and Savings Plans discount sharing for any accounts in that organisation, including the payer account.
    \item Renewal of Reserved Instances
    \item You can queue (schedule or reserve ahead of time) you reserved instances
    \item To renew a RI, just queue an RI purchase whenever the previous one expires
\end{itemize}